\documentclass[11pt,a4paper]{article}
\usepackage[utf8]{inputenc}
\usepackage{amsmath,amssymb,amsthm}
\usepackage{listings}
\usepackage[margin=2.5cm]{geometry}
\usepackage{xcolor}

\lstset{
  language=C++,
  basicstyle=\ttfamily\small,
  keywordstyle=\color{blue}\bfseries,
  commentstyle=\color{green!50!black},
  stringstyle=\color{red!70!black},
  numbers=left,
  numberstyle=\tiny\color{gray},
  breaklines=true,
  frame=single,
  tabsize=4,
  showstringspaces=false
}

\title{IOI 1993 -- Day 2, Task 2: The School Bus}
\author{Solution}
\date{}

\begin{document}
\maketitle

\section{Problem Statement}

A school bus starts at a depot, picks up students at $n$ locations, and
returns. The bus has capacity $C$, so it may need multiple round trips.
Each location $i$ has coordinates $(x_i, y_i)$ and $s_i$ students. Minimize
the total Euclidean distance traveled across all trips.

This is a variant of the \textbf{Capacitated Vehicle Routing Problem (CVRP)}.

\section{Solution}

\subsection{Two-Phase Bitmask DP}

For small $n$ (up to about 15), we use an exact algorithm in two phases.

\subsubsection{Phase 1: Optimal Tour for Each Subset (Held--Karp TSP)}

For each subset $S$ of locations whose total student count does not exceed
$C$, compute the shortest round trip from the depot visiting all locations
in $S$.

Define $\text{tsp}[S][v]$ as the minimum distance to visit every location in
$S$, starting from the depot and ending at location $v$. The recurrence is:
\[
  \text{tsp}[S][v] = \min_{u \in S \setminus \{v\}}
    \bigl(\text{tsp}[S \setminus \{v\}][u] + d(u, v)\bigr),
\]
with base case $\text{tsp}[\{v\}][v] = d(\text{depot}, v)$.

The round-trip cost for subset $S$ is:
\[
  \text{tour}[S] = \min_{v \in S}
    \bigl(\text{tsp}[S][v] + d(v, \text{depot})\bigr).
\]

\subsubsection{Phase 2: Set Partition DP}

Partition all $n$ locations into capacity-feasible trips to minimize total
cost:
\[
  \text{dp}[S] = \min_{\substack{T \subseteq S,\; T \ne \emptyset \\
    \text{students}(T) \le C}}
    \bigl(\text{tour}[T] + \text{dp}[S \setminus T]\bigr),
\]
with $\text{dp}[\emptyset] = 0$.

Subsets of $S$ are enumerated using the standard bitmask trick:
\texttt{for (T = S; T > 0; T = (T-1) \& S)}.

\section{Complexity Analysis}

\begin{itemize}
  \item \textbf{Phase 1 (TSP):} $O(2^n \cdot n^2)$ time, $O(2^n \cdot n)$
    space.
  \item \textbf{Phase 2 (partition):} $O(3^n)$ time (sum over all $S$ of the
    number of subsets of $S$ equals $3^n$).
  \item \textbf{Total:} $O(2^n \cdot n^2 + 3^n)$ time, $O(2^n \cdot n)$
    space.
  \item Feasible for $n \le 15$.
\end{itemize}

\section{Implementation}

\begin{lstlisting}
#include <iostream>
#include <vector>
#include <cmath>
#include <algorithm>
using namespace std;

const int MAXN = 16;
const double INF = 1e18;

int n, capacity;
int students[MAXN];
double dist_mat[MAXN+1][MAXN+1]; // 0 = depot, 1..n = locations

double tspDp[1 << MAXN][MAXN];
double tourCost[1 << MAXN];
int totalStudents[1 << MAXN];
bool feasible[1 << MAXN];
double partDp[1 << MAXN];

int main() {
    cin >> n >> capacity;

    double dx[MAXN+1], dy[MAXN+1];
    cin >> dx[0] >> dy[0]; // depot
    for (int i = 1; i <= n; i++)
        cin >> dx[i] >> dy[i] >> students[i-1];

    // Distance matrix
    for (int i = 0; i <= n; i++)
        for (int j = 0; j <= n; j++)
            dist_mat[i][j] = hypot(dx[i]-dx[j], dy[i]-dy[j]);

    int full = (1 << n) - 1;

    // Precompute total students per subset
    for (int S = 0; S <= full; S++) {
        totalStudents[S] = 0;
        for (int i = 0; i < n; i++)
            if (S & (1 << i))
                totalStudents[S] += students[i];
        feasible[S] = (totalStudents[S] <= capacity);
    }

    // Phase 1: TSP DP
    for (int S = 0; S <= full; S++)
        for (int i = 0; i < n; i++)
            tspDp[S][i] = INF;

    for (int i = 0; i < n; i++)
        tspDp[1 << i][i] = dist_mat[0][i+1];

    for (int S = 1; S <= full; S++) {
        if (!feasible[S]) continue;
        for (int v = 0; v < n; v++) {
            if (!(S & (1 << v))) continue;
            if (tspDp[S][v] >= INF) continue;
            for (int u = 0; u < n; u++) {
                if (S & (1 << u)) continue;
                int newS = S | (1 << u);
                if (!feasible[newS]) continue;
                double cost = tspDp[S][v] + dist_mat[v+1][u+1];
                if (cost < tspDp[newS][u])
                    tspDp[newS][u] = cost;
            }
        }
    }

    // Tour cost for each feasible subset
    for (int S = 0; S <= full; S++) {
        tourCost[S] = INF;
        if (!feasible[S] || S == 0) continue;
        for (int v = 0; v < n; v++) {
            if (!(S & (1 << v))) continue;
            double cost = tspDp[S][v] + dist_mat[v+1][0];
            tourCost[S] = min(tourCost[S], cost);
        }
    }

    // Phase 2: Set partition DP
    for (int S = 0; S <= full; S++)
        partDp[S] = INF;
    partDp[0] = 0;

    for (int S = 1; S <= full; S++) {
        for (int T = S; T > 0; T = (T - 1) & S) {
            if (feasible[T] && tourCost[T] < INF) {
                double val = partDp[S ^ T] + tourCost[T];
                if (val < partDp[S])
                    partDp[S] = val;
            }
        }
    }

    cout.precision(2);
    cout << fixed << partDp[full] << endl;

    return 0;
}
\end{lstlisting}

\section{Example}

\begin{verbatim}
Input:
4 6
0 0              (depot)
1 1 3            (location 1: 3 students)
3 1 2            (location 2: 2 students)
2 3 4            (location 3: 4 students)
4 2 1            (location 4: 1 student)

Total students = 10, capacity = 6.
At least 2 trips are needed.
\end{verbatim}

The DP finds the optimal partition into trips and the shortest tour for each.

\section{Notes}

\begin{itemize}
  \item This problem is NP-hard in general (it contains TSP as a special case
    when $C \ge \sum s_i$).
  \item The bitmask DP gives an exact solution for $n \le 15$.
  \item For larger $n$, heuristics (nearest-neighbor, savings algorithm) or
    metaheuristics (simulated annealing) are necessary.
  \item The subset enumeration identity $\sum_{S \subseteq [n]} 2^{|S|} = 3^n$
    explains the $O(3^n)$ time for Phase~2.
\end{itemize}

\end{document}
